\section*{Chapter \ref{ch:m1:exchange-traded-funds} --- Exchange-Traded Funds}

\begin{solution}{exo:m1:exchange-traded-funds:1}
$82.31/82.40 - 1 = -10.9$ basis points: a discount, outside the band of 9.
The \omterm{def:m1:exchange-traded-funds:ap}{authorised participant} buys ETF shares on the exchange, sells the basket
(short if need be), and redeems the ETF shares with the fund for the basket,
which closes its short. It earns about 1.9 basis points after costs.
\end{solution}

\begin{solution}{exo:m1:exchange-traded-funds:2}
The unit is worth $50\,000 \times 40 = \$2$ million; $500/2\,000\,000 = 2.5$
basis points. On ten units it is 0.25 basis point: the fee is fixed per
order, so the band is narrower for those who create in size, and small funds
with small creations have wider bands.
\end{solution}

\begin{solution}{exo:m1:exchange-traded-funds:3}
The fund returns $-2 \times -4\% = +8\%$; net assets become \$864 million.
$\Delta = \beta(\beta-1)Nr = 6 \times 800 \times (-0.04) = -\$192$ million:
the exposure must fall by 192 million. For an inverse fund ``less exposure''
means a larger short: it sells \$192 million of the index, after a fall. Check:
the short was 1\,600, it shrank with the index to 1\,536, and the target is $2
\times 864 = 1\,728$.
\end{solution}

\begin{solution}{exo:m1:exchange-traded-funds:4}
$2\times$: $1.20 \times 0.8182 - 1 = -1.82\%$. $3\times$: $1.30 \times 0.7273
- 1 = -5.45\%$. $-1\times$: $0.90 \times 1.0909 - 1 = -1.82\%$. The
three-times fund loses most: the loss grows like $\beta(\beta-1)$, which is
6 for $\beta = 3$ and 2 for both $\beta = 2$ and $\beta = -1$.
\end{solution}

\begin{solution}{exo:m1:exchange-traded-funds:5}
$\tfrac12\beta(\beta-1)\sigma^2 = 0.16$, $0.48$ and $0.48$: factors
$\mathrm{e}^{-0.16} = 0.852$ for $\beta = 2$ and $\mathrm{e}^{-0.48} = 0.619$
for both $\beta = 3$ and $\beta = -2$. A 10\% loss: $0.48\,T = -\ln 0.9 =
0.105$, so $T = 0.22$ year, about 55 trading days.
\end{solution}

\begin{solution}{exo:m1:exchange-traded-funds:6}
The gap between value and NAV is $6 \times 0.8^{k+1}$ on the $k$-th day after
the fall. NAV: 98.80, 97.84, 97.07, 96.46. Discount: 4.86\%, 3.92\%, 3.16\%,
2.55\%. It passes below 0.5\% on the eleventh day after the day of the fall.
A patient observer sees a discount of several percent ``closing'' over two
weeks while the ETF's price does not move at all.
\end{solution}

\begin{solution}{exo:m1:exchange-traded-funds:7}
With a band of 7 basis points an \omterm{def:m1:exchange-traded-funds:ap}{authorised participant} acts at 1\,638 of the
5\,000 steps and the mean absolute \omterm{def:m1:exchange-traded-funds:nav}{premium} is 3.28 basis points. With $c_B =
2$ the band is 5: 2\,219 actions and a mean absolute \omterm{def:m1:exchange-traded-funds:nav}{premium} of 2.42. The
investor who buys or sells at a random moment pays, on average, a smaller
deviation from fair value; cheaper underlying markets make tighter ETFs, and
more primary-market activity is the sign of a tighter product, not of a
stressed one.
\end{solution}

\begin{solution}{exo:m1:exchange-traded-funds:8}
The ``discount'' compares a live price with a NAV built from stale bond
valuations (\cref{prop:m1:exchange-traded-funds:stale}). It can close
entirely by the NAV falling to the ETF price, in which case the buyer of the
ETF earns nothing, or loses if the market continues down. To capture a real
discount one must buy the ETF \emph{and sell the basket} at the prices used
in the NAV, then redeem; in high-yield bonds in a falling market nobody will
buy the basket at those prices, which is exactly why the discount is there.
``Always closes'' is true and irrelevant; ``riskless'' is false.
\end{solution}

\begin{solution}{pb:m1:exchange-traded-funds:1}
\textbf{1.} $+3\%$ and $-3\%$: \$8.24 billion and \$2.91 billion.
\textbf{2.} $6 \times 8 \times 0.01 = \$0.48$ billion and $12 \times 3 \times
0.01 = \$0.36$ billion, both to buy: \$0.84 billion.
\textbf{3.} The last ten minutes trade \$30 billion: 2.8\%.
\textbf{4.} $\sigma_{10} = 1\% \times \sqrt{10/390} = 0.16\%$; impact $0.7
\times 0.16\% \times \sqrt{0.028} = 1.9$ basis points. Invisible.
\textbf{5.} $-15\%$ and $+15\%$: \$6.8 billion and \$3.45 billion.
\textbf{6.} $-\$2.4$ billion and $-\$1.8$ billion: \$4.2 billion to sell.
\textbf{7.} 14\%.
\textbf{8.} $\sigma_{10} = 0.48\%$; $0.7 \times 0.48\% \times \sqrt{0.14} =
12.6$ basis points.
\textbf{9.} The trade is $(6 \times 8 + 12 \times 3) = 84$ billion per unit of
return; a further $0.126\%$ adds \$0.11 billion, 2.5\% more. The feedback
exists and converges at once at this size; it would not if the funds were ten
times larger relative to the close.
\textbf{10.} $0.95 \times 1.04 \times 0.94 \times 1.05 \times 1.025 - 1 =
-0.05\%$: flat.
\textbf{11.} $0.85 \times 1.12 \times 0.82 \times 1.15 \times 1.075 - 1 =
-3.49\%$, against $-0.14\%$.
\textbf{12.} $1.15 \times 0.88 \times 1.18 \times 0.85 \times 0.925 - 1 =
-6.11\%$, against $+0.14\%$.
\textbf{13.} $-4.80\%$ in a week in which the index did nothing.
\textbf{14.} Both funds decay at the rate $\tfrac12\beta(\beta-1)\sigma^2$: 3
and 6 times the index's variance. The index exposures cancel; the decays add.
The pair is a short position in realised variance, paid for in five days of
4--6\% moves. (The mirror trade, short both funds, is long nothing and short
decay: it earns the variance but is exposed to a trend, in which one leg
compounds without limit, and to the cost of borrowing the shares.)
\textbf{15.} \omterm{def:m1:what-a-trading-firm-does:market-maker}{Market makers} and proprietary firms, who know the funds' assets
(published daily) and the day's return (public), hence the size and sign of
the trade, well before the close. They buy or sell ahead of it and provide
the other side in the closing auction (\cref{ch:m1:auctions}); their
competition is what keeps the impact near the estimate of question 8.
\textbf{16.} $\beta(\beta-1)$ is 12 against 6. After a rise the inverse
fund's assets fall while its short grows in value: both effects push its
leverage up, where for the long fund they partly offset.
\textbf{17.} $-33.3\%$ in one day. US market-wide circuit breakers halt
trading for the day when the S\&P 500 has fallen 20\%, which caps the one-day
loss of a fund on a broad US index below that level; funds on single stocks,
sectors or volatility have no such protection, and their prospectuses warn
that an investor can lose the entire investment in a single day.
\textbf{18.} The rebalancing trade as a fraction of closing volume: 14\% here,
2.8\% on the quiet day. It scales with assets times $\beta(\beta-1)$ times the
move, so it is largest exactly on the days when the close is already
strained.
\textbf{19.} \textbf{\$4.2 billion to sell.}
\textbf{20.} For someone who wants leverage for a day or a few days without
a margin account, and knows that beyond that horizon the product is a
position in the path, not only in the destination.
\end{solution}

\begin{solution}{iq:m1:exchange-traded-funds:1}
\omterm{def:m1:exchange-traded-funds:ap}{Authorised participants} can exchange the ETF's shares for the underlying
basket with the fund, in either direction, at NAV. If the ETF is dear they
sell it, buy the basket and create; if cheap, the reverse. The price stays
inside a band set by the cost of trading the basket, the ETF, the fee and
financing. Most of the time \omterm{def:m1:what-a-trading-firm-does:market-maker}{market makers} hedging with futures or the basket
do the work and creations are only the net residual.
\iqlookfor{both directions, the band rather than ``equal to NAV'', and the
fact that the fund itself does not trade.}
\end{solution}

\begin{solution}{iq:m1:exchange-traded-funds:2}
Almost never exactly. The fund delivers twice each \emph{daily} return; over
a year the result is the square of the index ratio times
$\exp(-\sigma^2 T)$. With 20\% volatility: $1.21 \times 0.961 = 1.163$, up
about 16\% before fees and financing; in a smooth trending year it can exceed
20\%.
\iqlookfor{the word ``daily'', a number, and the remark that trends help.}
\end{solution}

\begin{solution}{iq:m1:exchange-traded-funds:3}
No. The NAV uses Tokyo's close, ten hours old. Since then futures on the
Japanese index, the yen and US markets have moved; the ETF price is the
market's estimate of where Tokyo will open. Compare the ETF with a fair value
built from instruments that are open now. An arbitrage would require trading
the basket at the stale prices, which is impossible.
\iqlookfor{refusal of the word ``arbitrage'', and a concrete fair-value
construction.}
\end{solution}

\begin{solution}{iq:m1:exchange-traded-funds:4}
I am short \$50 million of the ETF. Immediately buy index futures for the
same delta: one trade, cheap, leaves basis risk between future and basket.
During the day other clients sell to me and reduce the position. Near the
close, for what remains, I either buy the basket (a program trade, often in
the closing auction so that my cost matches the NAV) and create ETF shares to
cover the short, selling the futures at the same time, or carry the short
ETF against futures overnight if borrowing the ETF is cheap and creation is
not worth the fee. My edge is the spread I charged minus hedge costs, fee and
financing.
\iqlookfor{hedge first, create last; awareness of the NAV strike at the
close and of inventory netting.}
\end{solution}

\begin{solution}{iq:m1:exchange-traded-funds:5}
Assets $N$, exposure $\beta N$. After return $r$: assets $N(1+\beta r)$,
exposure $\beta N(1+r)$, target $\beta N(1+\beta r)$; trade $\beta(\beta-1)Nr$.
The coefficient is positive for negative $\beta$ too: inverse funds buy after
rises and sell after falls, like leveraged ones, and a $-1\times$ fund trades
as much as a $2\times$ fund.
\iqlookfor{a clean three-line derivation and the sign result for inverse
funds.}
\end{solution}

\begin{solution}{iq:m1:exchange-traded-funds:6}
ETF mispriced: \omterm{def:m1:exchange-traded-funds:ap}{authorised participants}' balance sheets are constrained, so
the band has widened; forced sellers use the ETF because it is the only
liquid thing. NAV mispriced: bonds have not traded, valuations lag, and the
ETF leads. Evidence: prices of the bonds that did trade against their
valuations; the behaviour of NAV over the following days (does it fall to the
ETF?); credit-default-swap indices and futures as independent fair values;
redemption activity (if APs redeem heavily the discount is real to them); and
whether redemption baskets are being negotiated at prices below the
valuation.
\iqlookfor{both cases argued, and data that discriminates, especially the
subsequent convergence of NAV to the ETF.}
\end{solution}
